\section{Discussion}

\subsection{Interpretation of the Findings}

The results show that aggregate predictive performance is not sufficient to establish behavioral or explanatory stability in intrusion detection. Within the selected Rashomon set, models have comparable accuracy and macro-\gls{F1} scores, yet disagree on individual network events. This distinction matters because a model-selection procedure based only on a single aggregate score can select one member of a set of similarly plausible models without revealing where the members disagree. Our findings therefore instantiate predictive multiplicity \cite{marx2020predictive,hsu2022rashomon} in an intrusion-detection setting.

The disagreement is concentrated rather than uniform. Most test instances receive consistent predictions, whereas a smaller subset exhibits positive conflict ratios or substantial pairwise disagreement. These locally contested cases are precisely the instances for which an alert, triage decision, or downstream response may depend on an arbitrary choice among similarly performing models. Consequently, model multiplicity should not be interpreted as evidence that the detector is unusable; rather, it provides an additional uncertainty signal that is hidden by conventional performance reporting.

The explanation analyses lead to the same qualification. The models broadly identify similar destination-host traffic statistics as important, but their attribution magnitudes, feature rankings, and local \gls{SHAP} values vary. This is consistent with the observation that explanations of a single fitted model need not characterize all near-optimal models \cite{brunet2022implications,laberge2023partial}. In particular, the conflict-focused analysis shows that instability can affect features an analyst would plausibly use to understand a detection. Predictive conflict and explanation variability are related but not identical: a low-conflict sample may still receive different feature attributions, while prediction disagreement need not imply a large change in every attribution.

\subsection{Operational Implications for Cyber Defense}

For an intrusion-detection workflow, the practical implication is not to replace a detector with a committee by default. Instead, multiplicity measures can complement the usual model-selection and explanation steps. A deployment can report whether an event is stable across a validation-selected Rashomon set, use high-conflict events for analyst review or conservative escalation, and present feature attributions together with their variation across plausible models. This makes explicit what a single-model explanation otherwise obscures: whether the stated rationale is robust to a reasonable alternative model choice.

The approach is especially relevant where explanations support human decision-making. \gls{SHAP} is a useful common representation for comparing feature contributions \cite{lundberg2017unified}, but an explanation should not be interpreted as a unique causal account of an event. Reporting consensus features and attribution ranges can give analysts a more faithful picture of the evidence available from the model family. This complements existing calls for explainability in intrusion detection \cite{choras2020machine,szczepanski2020achieving} by adding a model-uncertainty perspective.

\subsection{Limitations and Threats to Validity}

The present study is limited to binary classification on \gls{NSLKDD} and a neural-network model family. The magnitude and structure of multiplicity may differ for modern traffic data, multiclass attacks, other feature representations, alternative classifiers, or operationally shifted traffic distributions. Moreover, the selected Rashomon set depends on the candidate-generation procedure, the performance metric, and the tolerance $\epsilon$. It should therefore be understood as an empirical approximation to the plausible model set, not as an exhaustive characterization of every equally good model.

There is also an important evaluation consideration. Rashomon-set membership should be determined using validation performance, while the test set remains reserved for reporting predictive and explanatory multiplicity. This avoids selecting models by the same test data used for the final analysis. The accompanying pipeline implements this separation and records the exact set membership; results produced by earlier test-selected artifacts should be treated as exploratory until they are regenerated through this validation-based workflow.

Finally, \gls{SHAP} values depend on the background data, the explained sample set, and the chosen aggregation metric. The reported variability measures therefore quantify explanation disagreement under a specified protocol rather than a model-independent ground truth. Future studies should evaluate sensitivity to these choices, quantify uncertainty across independent training origins, and test whether conflict- and explanation-variability signals improve analyst decisions or automated response policies.

